\documentclass[11pt]{article}
\usepackage[margin=.84in]{geometry}
\usepackage{amsmath,amssymb}
\usepackage[colorlinks=true,urlcolor=blue]{hyperref}
\setlength{\parindent}{0pt}
\setlength{\parskip}{.42em}
\setlength{\emergencystretch}{2em}
\newcommand{\dd}{\,dt}
\newcommand{\ch}{\chi_p}
\newcommand{\HT}{\href{https://jtnb.centre-mersenne.org/item/10.5802/jtnb.101.pdf}{Hildebrand--Tenenbaum, Theorem 1.6 and Corollary 1.7}}
\newcommand{\CM}{\href{https://arxiv.org/pdf/2404.08380v2}{Carneiro--Milinovich--Quesada-Herrera--Ramos, Theorems 1, 2, 4, 5 and the proof of Theorem 6}}
\newcommand{\DZZ}{\href{https://arxiv.org/pdf/2611.11389v1}{Dong--Zhao--Zhou, Theorem 1.4 and Sections 3--5 (v1)}}
\begin{document}
{\Large\bfseries Quadratic character sums over smooth integers at a full period}\\[3pt]
{\large The control at \(0.77(\log p)^2\) and the missing estimate at \(0.75(\log p)^2\)}

Let \(p\) run through all odd primes, \(\ch=(\cdot/p)\), \(L=\log p\),
\(\ell=\log L\), \(y=\eta L^2\) with fixed \(\eta>0\), and \(H=\log y\).
Write \(A(X,z)=\sum_{n\le X,\ P^+(n)\le z}\ch(n)\), \(A_p(y)=A(p,y)\),
\(\Psi(X,z)=\sum_{n\le X,\ P^+(n)\le z}1\), and \(P^+(1)=1\).
We assume Dirichlet GRH. All asymptotics are individual-prime estimates.
Put \(c_*=\mathcal C(1)^{-2}\) in the normalization of Carneiro et al.
The hand calculation below gives, for every fixed \(\eta>c_*\),
\begin{equation}
  |A_p(y)|\ \ll_\eta\ \Psi(p,y)\exp(-b_\eta L/H)
  \qquad (b_\eta>0).
\end{equation}
Their bounds \(0.7596<c_*<0.7615\) put \(0.77\) inside this result.
At \(0.75\) we isolate the remaining terms by size and frequency. We do
not obtain a nonresidue below this cutoff. \CM.

\section*{1. The size of one witness}
For large \(p\), \(y<p\), so all the signs in \(A_p(y)\) are \(\pm1\).
Thus \(A_p(y)=\Psi(p,y)\) if and only if \(n_p>y\). If \(q\le y\) is a
nonresidue prime, the pairs \((m,qm)\) for smooth \(m\le p/q\) have
opposite signs. Each negative member occurs at most twice, giving
\(\Psi(p,y)-A_p(y)\ge\Psi(p/q,y)\ge M\), where \(M=\Psi(p/y,y)\).
Using disjoint pairs strengthens this lower bound to
\begin{equation}
 \Psi(p,y)-A_p(y)\ \ge\
 2\sum_{j\ge1}\{\Psi(p/q^{2j-1},y)-\Psi(p/q^{2j},y)\}
 \ \ge\ (2-o_\eta(1))M .
\end{equation}
Pair each smooth integer of odd \(q\)-adic valuation with its quotient
by \(q\). The sum counts these disjoint pairs. For the last inequality,
if \(q\le y/9\) the original pairing already gives
\(\Psi(9p/y,y)\sim3M\). Otherwise the first difference alone is at least
\(M-\Psi(81p/y^2,y)=(1-o(1))M\).

For the scale in (2), let \(Z(s)=\prod_{v\le y}(1-v^{-s})^{-1}\), with
products over primes, and define \(\alpha\) by
\(\sum_{v\le y}\log v/(v^\alpha-1)=L\). Put \(Z=Z(\alpha)\).
The saddle-point formula and its local ratio estimate give
\begin{align}
 \alpha&=\tfrac12+(\log(2\sqrt\eta)+o_\eta(1))/H,\qquad
 \frac{M}{\Psi(p,y)}\sim y^{-\alpha}\sim\frac{1}{2\eta L},\label{eq:ratio}\\
 \log\Psi(p,y)&=\tfrac12 L+
    \big(1+\log(2\sqrt\eta)+o_\eta(1)\big)L/H .\label{eq:psi}
\end{align}
Indeed the PNT gives \(-(\log Z)'(\alpha)\sim y^{1-\alpha}/(1-\alpha)\),
\(\log Z\sim L/H\), and \(\phi_2=(\log Z)''(\alpha)\sim LH\). The first
estimate follows by bracketing \(\alpha=1/2+k/H\) for fixed \(k\).
The saddle formula gives \(p^\alpha Z/\Psi(p,y)\asymp\sqrt{LH}\).
Uniformly for \(1\le d\le y^K\), fixed \(K\), the local ratio is
\(\Psi(p/d,y)/\Psi(p,y)=(1+o(1))d^{-\alpha}\). Iterate Corollary 1.7
in steps of factor at most \(y\). The saddle changes by \(O_K(1/L)\)
since its derivative in \(\log p\), at fixed \(y\), is \(-1/\phi_2\).
\HT.

\newpage
\section*{2. What full-period cancellation removes}
Grouping by the unique largest prime and retaining its signed cofactor gives
\begin{equation}
 0=\sum_{n\le p}\ch(n)=A_p(y)+
       \sum_{y<r\le p\ {\rm prime}}\ch(r)A(p/r,r).
 \label{eq:buchstab}
\end{equation}
The cofactor may contain \(r\) again. The term \(r=p\) vanishes. There
is an independent estimate on a terminal part of this signed sum. For
\(1<U<\sqrt p\) with \(p/U>y\), write \(R\) for the sum over
\(y<r\le p/U\). Interchanging the prime and cofactor in the rest gives
\begin{equation}
 -A_p(y)-R=\sum_{m\le U}\ch(m)\!
     \sum_{p/U<r\le p/m\ {\rm prime}}\ch(r)
       \ \ll\ \sqrt{pU}\,L^2 .
 \label{eq:tail}
\end{equation}
Here \(r>\sqrt p>m\), so the smoothness condition on \(m\) is automatic.
The GRH bound \(\sum_{r\le t}\ch(r)\ll\sqrt t\log^2(pt)\) at both
endpoints and \(\sum_{m\le U}m^{-1/2}\ll\sqrt U\) suffice. In particular
write \(K_\eta=1+\log(2\sqrt\eta)\). If \(K_\eta>0\), take
\(U=\exp(\beta L/H)\) for any fixed \(0<\beta<2K_\eta\).
Equations (3)--(4) show the error in (6) is \(o_\eta(M)\) (with \(\beta\)
fixed). At \(0.75\), \(2K_\eta=2+\log3\), so one can take \(\beta=2\).
For \(\eta\le e^{-2}/4\) this comparison gives no positive \(\beta\).

Exceptionality at \(0.75\) would now force
\[
 R=-\Psi(p,y)+o(M).
\]
A uniform
\(R\ge-\Psi(p,y)+\delta M\), for any fixed \(\delta>0\), would suffice.
The first remaining prime term already has size \(M\). Let \(r\) be the
least prime above \(y\), so \(r/y\to1\). Under exceptionality,
\begin{equation}
 \ch(r)A(p/r,r)=\sum_{k\ge1}\ch(r)^k\Psi(p/r^k,y)
              =(\ch(r)+o(1))M .
\end{equation}
The first term uses (3)'s local ratio. Those with \(k=2,3\) are \(o(M)\),
and the rest are at most \((L/\log r)\Psi(p/r^4,y)=o(M)\) in absolute
value. The unknown sign can go either way in this decomposition. A joint
estimate could still control \(R\). Removing the first prime term by its
absolute value alone costs the full one-witness scale.

The same grouping prevents a direct transfer from \(z=0.77L^2\) to
\(y=0.75L^2\). Its exact cost is
\(A(p,z)-A(p,y)=\sum_{y<r\le z}\ch(r)A(p/r,r)\), over primes.
Even a very small \(A(p,z)\) leaves this signed difference.
The unsigned counts themselves have ratio
\(\Psi(p,z)/\Psi(p,y)=\exp\{(\tfrac12\log(0.77/0.75)+o(1))L/H\}\)
by (4).

\textbf{Where the November estimate loses the cutoff.} At \(x=p\) the
preprint takes \(y=(L+L^2)\ell^C\). On
\(\sigma_b=1/2+b\log\ell/\ell\), Lemmas 3.1--3.2 give an upper logarithm
\(L/2+bW+O(L\ell^{2-2b}+\ell)\), \(W=L\log\ell/\ell\).
The error comes from the log-Euler-product approximation (3.3), with its
factor \((\log v)^2 v^{1/2-\sigma_b}\). Lemma 4.1 gives a lower logarithm
\(L/2+(C/4+o(1))W\). Section 5 chooses \(3/2<b<C/4-\gamma\).
At a fixed multiple of \(L^2\), (4) puts the true surplus at order at
most \(L/\ell=o(W)\). Even an exact unsigned count cannot pay for the
contour displacement \(bW\). Setting the ordinary character sum to zero
does not improve the separate bound on the finite product. \DZZ.

\newpage
\section*{3. The Fourier input at \(0.77\)}
For each fixed \(\eta>c_*\) there are constants \(d,\delta\in(0,1)\),
depending on \(\eta\), such that for each sign \(\varepsilon=\pm1\),
\begin{equation}
 \#\{d y\le v\le y:\ v\ {\rm prime},\ \ch(v)=\varepsilon\}
       \ge\delta\,y/H .
 \label{eq:both}
\end{equation}
We give the quantitative deduction because merely finding one prime of
each sign is insufficient for (1). Take an admissible fixed test \(F\)
with \(\widehat F\) real, smooth and supported in \([-N,N]\). Put
\[
 D(F)=\int_{-\infty}^0\widehat F(t)e^{\pi t}\,dt
              -\int_0^\infty|\widehat F(t)|e^{\pi t}\,dt.
\]
By the extremal characterization, \(\inf_{D(F)>0}\|F\|_1/(2\pi D(F))
=\sqrt{c_*}\). Choose \(F\) with
\(\sqrt\eta D(F)>\|F\|_1/(2\pi)\), let \(\Delta=H/(2\pi)\) and
\(\widehat h(t)=(\widehat F(t-\Delta)+\widehat F(-t-\Delta))/2\).
The explicit formula, the zero count and the gamma estimate in the proof
of Theorem 6 give
\[
 \left|\pi^{-1}\sum_n\Lambda(n)\ch(n)n^{-1/2}
                 \widehat h(\log n/(2\pi))\right|
 \le L\|F\|_1/(2\pi)+o_F(L).
\]
Here \(\widehat h(0)=0\). Both signs of this inequality follow from the
absolute zero majorant. The supported positive integers have
\(d y\le n\le y/d\) for \(d=e^{-2\pi N}\). Proper prime powers
contribute \(O_F(\log^2 y)\).

For either \(\varepsilon\), use prime coefficient \(-\varepsilon\ch(v)\).
Putting \(+1\) below \(y\), and minus the absolute value of the weight
above \(y\), gives \(\sqrt y D(F)+o_F(\sqrt y)\) by the PNT and the
change of variable \(v=y e^{2\pi t}\). Restoring each coefficient below
\(y\) with \(\ch(v)=\varepsilon\) costs at most \(O_F(H/\sqrt y)\).
The strict margin in the choice of \(F\) proves (8). At \(\eta=0.75\),
the lower bound \(c_*>0.7596\) prevents this choice for every admissible
test. This is the specific arithmetic input lost on moving the cutoff.
\CM.

\textbf{A direct count check on the sign input.} Match distinct positive
primes in (8) with distinct negative primes to form at least
\(\lfloor\delta y/H\rfloor\) disjoint prime pairs. For each pair \((r,q)\)
and smooth \(m\le p/y\), exactly one of \(rm,qm\) has negative sign.
A given negative integer can occur at most \(L/\log(dy)\) times since
each occurrence uses a different prime divisor at least \(d y\).
Thus
\[
 \Psi(p,y)-A_p(y)\ \ge\
 2\lfloor\delta y/H\rfloor\,M\,\log(dy)/L \gg_\eta\Psi(p,y).
\]
Interchanging the sign being counted gives the same lower bound for
\(\Psi(p,y)+A_p(y)\). This gives a fixed fractional gap already.
The next two steps show how the same prime counts give nearly equal
numbers of the two signs.

\newpage
\section*{4. Uniform damping away from frequency zero}
For \(f=\ch\) or \(f=1\) put
\[
 F_y(s;f)=\prod_{v\le y}(1-f(v)v^{-s})^{-1},\qquad
 U_f(t)=|F_y(\alpha+it;f)|/Z .
\]
The Euler expansion for the logarithm gives
\[
 -\log U_f(t)\ \ge\ \sum_{v\le y}
       \frac{1-f(v)\cos(t\log v)}{v^\alpha}.
\]
All prime-power losses omitted here are nonnegative. For \(\eta>c_*\)
choose \(\tau>0\) with \(\tau|\log d|\le1/4\), using \(d\) from (8).
For \(|t|\le\tau\), \(\cos(t\log v)\) differs from \(\cos(t\log y)\)
by at most \(1/4\) on \([dy,y]\). Whichever sign the latter cosine has,
one of the two sets in (8) has loss at least \(3/4\) per prime before
the weight \(v^{-\alpha}\). Hence \(U_{\ch}(t)\le e^{-c_\eta L/H}\)
on this whole interval.

For \emph{every} fixed \(\eta>0\), including \(0.75\), use
\(1-f(v)\cos z\ge 1-|\cos z|\ge(1-\cos(2z))/4\).
Fix a nonzero nonnegative smooth \(g\) compactly supported in \((1/2,1)\)
and write \(G(s)=\int_0^\infty g(u)u^{s-1}\,du\). Under RH for \(\zeta\),
\begin{equation}
 \sum_{v\ {\rm prime}}(\log v)g(v/y)v^{iu}
 =y^{1+iu}G(1+iu)+O_g\big(\sqrt y[\log(2+|u|)+\log^2 y]\big).
 \label{eq:twist}
\end{equation}
Indeed Mellin inversion for the von Mangoldt sum shifts
\(-(\zeta'/\zeta)(s)y^{s+iu}G(s+iu)\) to \(\Re s=-1\).
The zeros contribute \(O_g(\sqrt y\log(2+|u|))\) by rapid decay and
\(O(\log(2+|v_0|))\) zeros per unit interval at height \(v_0\).
The new vertical integral is \(O_g(y^{-1}\log(2+|u|))\) by the functional
equation. Removing proper powers costs \(O_g(\sqrt y\log^2 y)\).
\href{https://people.maths.ox.ac.uk/greenbj/papers/primenumbers.pdf}{Green, chapter 7, smoothed explicit formula}.

For any fixed \(\tau>0\),
\(\sup_{|u|\ge2\tau}|G(1+iu)|<G(1)\) by strict triangle inequality
and decay. Near zero, \(G(1)-|G(1+iu)|\gg_g u^2\) by the positive
variance of \(\log v\) under the normalized measure \(g(v)\,dv\).
Subtract the real part of (9) at \(u=2t\) from its value at \(0\).
For \(\tau\le |t|\le\exp(L/H)\) the error is \(o_\eta(y)\).
For \(1/H\le|t|\le\tau\) it is \(o_\eta(y/H^2)\).
Since \(g\) is bounded and \(v^{-\alpha}/\log v\gg_\eta 1/(LH)\)
on its support after scaling by \(y\), this proves, for both \(f\),
\begin{align}
 U_f(t)&\le e^{-c_{\eta,\tau}L/H}
           &&(\tau\le |t|\le e^{L/H}),\label{eq:high}\\
 U_f(t)&\le e^{-c_{\eta,\tau}L/H^3}
           &&(1/H\le |t|\le\tau).\label{eq:middle}
\end{align}
For (11) the lower bound for \(G(1)-|G(1+iu)|\) on compact ranges uses
strictness away from zero as well. Estimate (10) uses only the reality
of the prime signs and RH for \(\zeta\). Under (8) the stronger rate
\(e^{-c_\eta L/H}\) holds for \(f=\ch\) down to \(t=0\).

\newpage
\section*{5. Returning to the sharp endpoint}
Choose \(T=\exp(aL/H)\) with \(0<a<1\) sufficiently small compared with
the constant in (10), and put \(\epsilon=T^{-1/2}\). The cutoff average
\[
 \frac1{2\epsilon}\int_{-\epsilon}^{\epsilon}
       \sum_{\substack{n\le pe^v\\P^+(n)\le y}}f(n)\,dv
 =\frac1{2\pi}\int_{\mathbb R}
   \frac{p^sF_y(s;f)}s\,\frac{\sinh(\epsilon s)}{\epsilon s}\,dt,
 \qquad s=\alpha+it ,
\]
follows by integrating Perron's kernel. The smoothed kernel is absolutely
integrable and the finite Euler product has an absolutely convergent
Dirichlet series on this line. The integral for \(|t|>T\) costs
\(O(p^\alpha Z/(\epsilon T))\).

To remove the averaging, use the probability measure on points \(\log n\)
with mass \(n^{-\alpha}/Z\) at each \(y\)-smooth \(n\). Its Fourier
transform in absolute value is \(U_1(t)\). Its mass within \(\epsilon\)
of any point is
\[
 \ll\epsilon\int_{|t|\le1/\epsilon}U_1(t)\dd\ll_{\eta,\tau}\epsilon .
\]
For the first inequality majorize the interval indicator by a constant
times \((\sin(v/(2\epsilon))/(v/(2\epsilon)))^2\) in the centered log
variable \(v\), then use its triangular Fourier transform. The second
inequality uses \(U_1\le1\) up to fixed \(\tau\) and (10) afterwards,
with \(a\) chosen as above. Thus the difference between the averaged and
sharp sums is \(O_{\eta,\tau}(\epsilon p^\alpha Z)\), for either sign
choice, uniformly at the endpoint.

For \(\eta>c_*\) apply the strong bound for \(U_{\ch}(t)\) on
\([-T,T]\). The remaining kernel has integral in absolute value
\(O(\log(2T))\) after removing \(F_y\) and \(p^\alpha\).
The tail and smoothing errors are exponentially small as well.
Since \(p^\alpha Z/\Psi(p,y)\asymp\sqrt{LH}\), reducing the positive
constant in the exponent absorbs the polynomial factors. This proves (1).

For any fixed \(\eta>0\), (10)--(11) instead yield the low-frequency
formula, with \(f=\ch\) or \(1\),
\begin{equation}
 \sum_{n\le p,\ P^+(n)\le y} f(n)=
 \frac{p^\alpha}{2\pi}\int_{-1/H}^{1/H}
       \frac{F_y(\alpha+it;f)p^{it}}{\alpha+it}\dd
 +O_\eta\!\left(\Psi(p,y)e^{-c_\eta L/H^3}\right).
\end{equation}
On this short interval the smoothing factor is \(1+O(\epsilon^2)\).
The error is \(o(M)\). At \(0.75\), subtracting the two copies of (12)
therefore places the entire one-witness question in the displayed
integral. If every small prime is a residue the two products there
coincide identically, including at \(t=0\).

One can bound the central integral at any fixed \(\eta\) in terms of the
actual negative prime mass \(D_p=\sum_{v\le y,\ \ch(v)=-1}v^{-\alpha}\).
On \(|t|\le1/H\), the first-power loss is at least
\(2\cos(1)D_p+c_\eta LHt^2\). Separate \(1-\cos(t\log v)\) from the
additional \(2\cos(t\log v)\) at negative primes, and sum the former
over \([y/2,y]\) by the PNT. Gaussian integration in (12) yields
\[
 \frac{|A_p(y)|}{\Psi(p,y)}
 \ll_\eta e^{-2\cos(1)D_p}+e^{-c_\eta L/H^3}.
\]
At \(0.75\) this estimate still allows \(D_p=0\).

\textbf{Consequence for \(n_p\).} At \(0.77\), (1) puts each sign on
\((1/2+o(1))\Psi(p,y)\) smooth integers. At \(0.75\), (6) controls the
largest-prime tail and (12) controls the outer frequencies beyond the
precision in (2). Neither provides a strict gap for the remaining signed
sum. The growth-rate bound on \(n_p\) is unchanged. These are hand
derivations. Priority for (1) has not been checked.
\end{document}
